\documentclass[../../main.tex]{subfiles}
\begin{document}

{\bf [解答]}
\paragraph{(1)}

\begin{figure}[htb]
  \centering
  \begin{tikzpicture}[scale=1.4, >=stealth, every node/.style={font=\small}]
    % 頂点座標
    \coordinate (A) at (0, 3.0);
    \coordinate (B) at (-2.0, -0.2);
    \coordinate (C) at (-0.3, -1.0);
    \coordinate (D) at (2.0, -0.2);

    % 底面 BCD の重心 H
    \coordinate (H) at (-0.1, -0.47);

    % 奥の稜線 BD
    \draw[thick, dashed, gray!70!black] (B) -- (D);

    % 底面の稜線 BC, CD
    \draw[line width=1.1pt, gray!85!black] (B) -- (C) -- (D);

    % 垂線 AH
    \draw[densely dotted, line width=1.0pt, gray!70!black] (A) -- (H);
    \node[circle, fill=gray!80!black, inner sep=1.1pt] at (H) {};
    \node[below=2pt, font=\footnotesize] at (H) {$H$};

    % AP = AQ = AR = x (比率 0.6)
    \def\rRatio{0.6}
    \coordinate (P) at ($ (A)!\rRatio!(B) $);
    \coordinate (Q) at ($ (A)!\rRatio!(C) $);
    \coordinate (R) at ($ (A)!\rRatio!(D) $);

    % 垂足 P', Q', R' (P - P' は AH と平行)
    \coordinate (Pp) at ($ (H)!\rRatio!(B) $);
    \coordinate (Qp) at ($ (H)!\rRatio!(C) $);
    \coordinate (Rp) at ($ (H)!\rRatio!(D) $);

    % 三角柱 PQR の上面ハッチング
    \fill[pattern=north east lines, pattern color=blue!45] (P) -- (Q) -- (R) -- cycle;

    % 柱の側面投影線
    \draw[thick, dashed, blue!70!black] (P) -- (Pp);
    \draw[line width=1.1pt, blue!85!black] (Q) -- (Qp);
    \draw[line width=1.1pt, blue!85!black] (R) -- (Rp);

    % 底面 P'Q'R' の稜線
    \draw[thick, dashed, blue!70!black] (Pp) -- (Rp);
    \draw[line width=1.1pt, blue!85!black] (Pp) -- (Qp) -- (Rp);

    % 手前の稜線 AB, AC, AD
    \draw[line width=1.1pt, gray!85!black] (A) -- (B);
    \draw[line width=1.1pt, gray!85!black] (A) -- (C);
    \draw[line width=1.1pt, gray!85!black] (A) -- (D);

    % 三角柱上面 PQR
    \draw[line width=1.2pt, blue!85!black] (P) -- (Q) -- (R) -- cycle;

    % 各点ノード
    \node[above=2pt] at (A) {$A$};
    \node[left=2pt] at (B) {$B$};
    \node[below left=1pt] at (C) {$C$};
    \node[right=2pt] at (D) {$D$};

    \node[circle, fill=blue!85!black, inner sep=1.2pt] at (P) {};
    \node[left=2pt, blue!85!black] at (P) {$P$};
    \node[circle, fill=blue!85!black, inner sep=1.2pt] at (Q) {};
    \node[left=2pt, blue!85!black] at (Q) {$Q$};
    \node[circle, fill=blue!85!black, inner sep=1.2pt] at (R) {};
    \node[right=2pt, blue!85!black] at (R) {$R$};

    \node[circle, fill=blue!70!black, inner sep=1.0pt] at (Pp) {};
    \node[below left=-1pt, font=\footnotesize, blue!70!black] at (Pp) {$P'$};
    \node[circle, fill=blue!85!black, inner sep=1.0pt] at (Qp) {};
    \node[below=2pt, font=\footnotesize, blue!85!black] at (Qp) {$Q'$};
    \node[circle, fill=blue!85!black, inner sep=1.0pt] at (Rp) {};
    \node[below right=-1pt, font=\footnotesize, blue!85!black] at (Rp) {$R'$};

    % 寸法 x の表示
    \draw[<->, font=\footnotesize] ($ (A) + (-0.25, 0) $) -- ($ (P) + (-0.25, 0) $)
    node[midway, left] {$x$};
  \end{tikzpicture}
  \caption{正四面体 $ABCD$ と内接する三角柱 $PQR-P'Q'R'$}
  \label{fig:tetrahedron_prism}
\end{figure}

\begin{figure}[htb]
  \centering
  \begin{tikzpicture}[>=stealth, every node/.style={font=\small}]
    % --- Panel 1: 切断面 ABH における相似 ---
    \begin{scope}[shift={(0,0)}]
      \coordinate (H) at (0, 0);
      \coordinate (A) at (0, 2.6);
      \coordinate (B) at (1.84, 0);

      % 比率 r = x/a = 0.6
      \def\rRatio{0.6}
      \coordinate (P) at ($ (A)!\rRatio!(B) $);
      \coordinate (Pp) at ($ (H)!\rRatio!(B) $);

      % 直角記号
      \draw[gray!70!black] (0, 0.2) -- (0.2, 0.2) -- (0.2, 0);
      \draw[gray!70!black] ($ (Pp) + (0, 0.2) $) -- ($ (Pp) + (0.2, 0.2) $) -- ($ (Pp) + (0.2, 0) $);

      % 三角形 ABH の辺
      \draw[line width=1.1pt, gray!85!black] (A) -- (H) -- (B) -- cycle;

      % 垂線 PP'
      \draw[line width=1.2pt, blue!85!black] (P) -- (Pp);

      % 点
      \node[circle, fill, inner sep=1.2pt] at (A) {};
      \node[above=2pt] at (A) {$A$};
      \node[circle, fill, inner sep=1.2pt] at (H) {};
      \node[below left=1pt] at (H) {$H$};
      \node[circle, fill, inner sep=1.2pt] at (B) {};
      \node[below right=1pt] at (B) {$B$};

      \node[circle, fill=blue!85!black, inner sep=1.3pt] at (P) {};
      \node[above right=1pt, blue!85!black] at (P) {$P$};
      \node[circle, fill=blue!85!black, inner sep=1.3pt] at (Pp) {};
      \node[below=2pt, blue!85!black] at (Pp) {$P'$};

      % 寸法表示
      \node[left=2pt] at ($ (A)!0.5!(H) $) {$|AH|$};
      \node[right=2pt, blue!85!black] at ($ (P)!0.5!(Pp) $) {$|PP'|$};
      \node[above right=1pt] at ($ (A)!0.5!(P) $) {$x$};
      \node[above right=1pt] at ($ (P)!0.5!(B) $) {$a-x$};

      \node[below=0.8cm] at (0.9, 0) {(a) 切断面 $\triangle ABH$（$\triangle BPP' \sim \triangle BAH$）};
    \end{scope}

    % --- Panel 2: 側面正三角形 ABC と PQ ---
    \begin{scope}[shift={(4.5, 0)}]
      \coordinate (B2) at (0, 0);
      \coordinate (C2) at (2.6, 0);
      \coordinate (A2) at (1.3, 2.25);

      \def\rRatio{0.6}
      \coordinate (P2) at ($ (A2)!\rRatio!(B2) $);
      \coordinate (Q2) at ($ (A2)!\rRatio!(C2) $);

      % 三角形 APQ の薄いハッチング
      \fill[pattern=north east lines, pattern color=blue!30]
      (A2) -- (P2) -- (Q2) -- cycle;

      % 正三角形 ABC
      \draw[line width=1.1pt, gray!85!black] (A2) -- (B2) -- (C2) -- cycle;

      % 線分 PQ
      \draw[line width=1.2pt, blue!85!black] (P2) -- (Q2);

      % 点
      \node[circle, fill, inner sep=1.2pt] at (A2) {};
      \node[above=2pt] at (A2) {$A$};
      \node[circle, fill, inner sep=1.2pt] at (B2) {};
      \node[below left=1pt] at (B2) {$B$};
      \node[circle, fill, inner sep=1.2pt] at (C2) {};
      \node[below right=1pt] at (C2) {$C$};

      \node[circle, fill=blue!85!black, inner sep=1.3pt] at (P2) {};
      \node[left=2pt, blue!85!black] at (P2) {$P$};
      \node[circle, fill=blue!85!black, inner sep=1.3pt] at (Q2) {};
      \node[right=2pt, blue!85!black] at (Q2) {$Q$};

      % 寸法表示
      \node[left=1pt] at ($ (A2)!0.5!(P2) $) {$x$};
      \node[right=1pt] at ($ (A2)!0.5!(Q2) $) {$x$};
      \node[below=1pt, blue!85!black] at ($ (P2)!0.5!(Q2) $) {$x$};
      \node[below=2pt] at ($ (B2)!0.5!(C2) $) {$a$};

      \node[below=0.8cm] at (1.3, 0) {(b) 正三角形 $\triangle ABC$ と線分 $PQ$};
    \end{scope}
  \end{tikzpicture}
  \caption{切断面における相似と正三角形の各部}
  \label{fig:cross_section_and_equilateral}
\end{figure}

$A$から平面$BCD$に下ろした垂足$H$とする．$ABCD$が正四面体だから
$|AH|=\dfrac{\sqrt{6}}{3}a$である．
$|AP|=x \ (0<x<a)$とおく．すると相似から
\begin{align}
  |PP'|&=\dfrac{a-x}{a}|AH|=\dfrac{\sqrt{6}}{3}(a-x) \\
  |PQ|&=x
\end{align}
である．$\triangle PQR$は正三角形でその面積
\begin{align}
  S(x)=\dfrac{\sqrt{3}}{4}|PQ|^2=\dfrac{\sqrt{3}}{4}x^2
\end{align}
だから，三角柱の体積$V(x)$は
\begin{align}
  V(x)&=(\triangle PQR)|PP'| \\
  &=\dfrac{\sqrt{6}}{3}(a-x)S(x) \\
  &=\dfrac{\sqrt{2}}{4}x^2(a-x)
\end{align}
となる．$x,a-x>0$からAM-GMより
\begin{align}
  V(x)
  &=\sqrt{2}\left(\dfrac{x}{2}\right)^2(a-x) \nonumber\\
  &\le\sqrt{2}\left(\dfrac{a}{3}\right)^3\label{eq:1}
\end{align}
である．等号成立は$\dfrac{x}{2}=a-x\Longleftrightarrow x=\dfrac{2}{3}a$である．(これは$0<x<a$を満たす．)
以上から求める値は$|AP|=\dfrac{2}{3}a$である．$\cdots$(答)

\paragraph{(2)}
\cref{eq:1}から$V_0=\dfrac{\sqrt{2}}{27}a^3$であり，また
\begin{align}
  V
  =\dfrac{1}{3}S\left(x=\dfrac{2}{3}a\right)|AH|
  =\dfrac{\sqrt{2}}{12}a^3
\end{align}
だから，
\begin{align}
  \dfrac{V_0}{V}=\dfrac{\dfrac{\sqrt{2}}{27}a^3}{\dfrac{\sqrt{2}}{12}a^3}=\dfrac{4}{9}
\end{align}
である．$\cdots$(答)
\newpage
\end{document}
